\documentclass[12pt]{article}
\usepackage{amsfonts,amsthm,amsmath,amssymb,mathtools}
\usepackage{mathrsfs}
\usepackage{enumitem}
\usepackage{tikz-cd}
\usepackage{geometry}
\usepackage{mlmodern}

\geometry{letterpaper, portrait, margin=1in}

\setcounter{section}{0}

\renewcommand{\thesection}{\S\arabic{section}}
\renewcommand{\thesubsection}{\arabic{section}}
\newcommand{\x}{\mathbf{x}}

\newtheorem{thm}{Theorem}
\newenvironment{theorem}[1]{
	\renewcommand{\thethm}{#1}
	\begin{thm}
		}{\end{thm}}

\newtheorem{lma}{Lemma}
\newenvironment{lemma}[1]{
	\renewcommand{\thelma}{#1}
	\begin{lma}
		}{\end{lma}}

\theoremstyle{definition}
\newtheorem{ex}{}
\newenvironment{exercise}[1]{
	\renewcommand{\theex}{#1}
	\begin{ex}
		}{\end{ex}}

\title{Homework 7}
\author{Connor Mooneyhan}
\date{March 16, 2026}

\begin{document}

\maketitle

\begin{ex}
	A semi-norm on a vector space $X$ is a map $g : X \to \mathbb{R}$ that is always non-negative, homogeneous, and satisfies the triangle inequality.
	Show that if $p$ is a semi-norm then $p(0) = 0$ and $|p(u) - p(v)| \leq p(u-v)$.
	Hence if $p(u) = 0$ implies $u = 0$ (assuming that's what was meant and the problem statement was initially typed incorrectly as $p = 0$) then $p$ is a norm.
\end{ex}
\begin{proof}
	First, if $p$ is a semi-norm, then \begin{equation*}
		p(0) = p(0 \cdot 0) = 0 \cdot p(0) = 0.
	\end{equation*}
	Also, for any $u, v \in X$, \begin{align*}
		         & p(u) = p((u - v) + v) \leq p(u - v) + p(v) \\
		\implies & p(u) - p(v) \leq p(u - v).
	\end{align*}
	Swapping the roles of $u$ and $v$, we get \begin{align*}
		         & p(v) - p(u) \leq p(u - v) = p(-1 \cdot (v - u)) = p(v - u) \\
		\implies & p(u) - p(v) \geq -p(v-u) = -p(u - v),
	\end{align*}
	where the second line comes from multiplying the first by $-1$.
	Combining the final inequalities, we get \begin{equation*}
		|p(u) - p(v)| \leq p(u - v)
	\end{equation*}
	as desired.

	Finally, if $p(u) = 0$ implies $u = 0$, then since $p$ already satisfies positivity, homogeneity, and the triangle inequality, the only condition left for it to be a norm was definiteness, and $p(u) = 0 \implies u = 0$ is exactly that property.
	Thus $p$ is a norm.
\end{proof}

\begin{ex}
	Show that if $p$ is a semi-norm, then the set \begin{equation*}
		N \coloneq \left\lbrace u : p(u) = 0 \right\rbrace
	\end{equation*}
	is a subspace of $X$, and one can define a norm on the quotient space $X/N$ by $||[u]|| = p(u)$.
\end{ex}
\begin{proof}
	To show that $N$ is a subspace, we just need to show that it is closed under the operations.
	If $u, v \in N$, then \begin{equation*}
		0 \leq p(u + v) \leq p(u) + p(v) = 0 \implies p(u+v) = 0,
	\end{equation*}
	where the first inequality comes from nonnegativity and the second comes from the triangle inequality.
	Thus $u + v \in N$.
	Now if $c \in \mathbb{R}$ (wlog) and $u \in N$, then \begin{equation*}
		p(c\cdot u) = |c|\cdot p(u) = |c|\cdot 0 = 0,
	\end{equation*}
	so $c\cdot u \in N$.

	First we show that $||\cdot||$ is well defined on $X/N$.
	If $[u] = [v]$, then $u - v \in N$, so $p(u - v) = 0$.
	Using the fact that \begin{equation*}
		p(u) = p(u - v + v) \leq p(u - v) + p(v) \implies p(u) - p(v) \leq p(u - v),
	\end{equation*}
	we have $p(u) - p(v) \leq p(u - v) = 0$, so $p(u) = p(v)$.
	Thus $||\cdot||$ is well defined.

	Now we show that $||\cdot||$ is a norm on $X/N$.
	Nonnegativity comes from the fact that $||[u]|| = p(u) \geq 0$ for all $[u] \in X/N$.
	Similarly, homogeneity follows by \begin{equation*}
		||[cu]|| = p(cu) = |c| p(u) = |c|\cdot ||[u]||.
	\end{equation*}
	The triangle inequality as well is immediate: \begin{align*}
		||[u + v]|| & = p(u + v)           \\
		            & \leq p(u) + p(v)     \\
		            & = ||[u]|| + ||[v]||.
	\end{align*}
	So we finally get to definiteness, where we must first show that $0 \in N$.
	Indeed, by homogeneity, \begin{equation*}
		p(0) = p(0 \cdot 0) = 0 \cdot p(0) = 0.
	\end{equation*}
	Now if $[u] \in X/N$ such that $||[u]|| = 0$, then $p(u) = 0$, so $u \in N$, and thus $[u] = [0]$.
	This gives us definiteness, so $|| \cdot ||$ is a norm on $X/N$.
\end{proof}

\begin{ex}
	Construct a sequence of functions $(f_n)$ in $C[0, 1]$ such that $(f_n) \in [0] \in L^p[0, 1]$ (here $[0]$ is the additive identity of $L^p[0,1]$) for every $p \geq 1$, but $(f_n(x))$ diverges for every $x \in [0, 1]$.
\end{ex}
\begin{proof}
	We will take the suggested ``trapezoidal'' route by describing each $f_n$ in the ``$k$'' subset $S_k$ of $(f_n)$, namely the trapezoids that cover $[0, 1/2^k], [1/2^k, 2/2^k], \dots, [(2^k-1)/2^k, 1]$.
	For the first $f_n$ in $S_k$, we have the formula \begin{equation*}
		f_n(x) = \begin{cases}
			1         & x \in \left[0, \frac{1}{2^k} \right]                        \\
			-2^kx + 2 & x \in \left[\frac{1}{2^k} \leq x \leq \frac{2}{2^k} \right] \\
			0         & x \in \left[\frac{2}{2^k}, 1 \right].
		\end{cases}
	\end{equation*}
	For the final $f_n$ in $S^k$, we have the formula \begin{equation*}
		f_n(x) = \begin{cases}
			0              & x \in \left[0, \frac{2^k-2}{2^k} \right]                      \\
			2^kx + 2 - 2^k & x \in \left[ \frac{2^k - 2}{2^k}, \frac{2^k - 1}{2^k} \right] \\
			1              & x \in \left[\frac{2^k-1}{2^k}, 1 \right].
		\end{cases}
	\end{equation*}
	Finally, for the ``middle'' $f_n$'s in $S^k$, for some integer $m$ where $1 \leq m \leq 2^k-2$ we have the formula \begin{equation*}
		f_n(x) = \begin{cases}
			0 & x \in \left[ 0, \frac{m-1}{2^k} \right]                  \\
			2^kx+1-m & x \in \left[ \frac{m-1}{2^k}, \frac{m }{2^k} \right]                                  \\
			1 & x \in \left[ \frac{m }{2^k}, \frac{m + 1}{2^k} \right] \\
			-2^kx+m+2 & x \in \left[ \frac{m+1}{2^k}, \frac{m+2}{2^k} \right]                                  \\
			0 & x \in \left[ \frac{m+2}{2^k}, 1 \right].
		\end{cases}
	\end{equation*}
  Indeed, this converges pointwise for no $x$ since the sequence $(f_n(x))$ for a fixed $x$ is an unending chain of zeroes, ones, and some intermediary numbers for the ``edges'' of the trapezoids.
  The important part is that the zeroes and ones keep occuring in the sequence, so it cannot converge.
  However, the integral of these functions continues to get smaller and smaller, and in fact limits to 0, since the support of each $f_n^p$ is at most $3/2^k$ and the height of $f_n^p$, so \begin{equation*}
    \left( \int_0^1 f_n^p \right)^{1/p} \leq \frac{3}{2^k}.
  \end{equation*}
  Thus, $(f_n) \in L^p[0,1]$ but does not converge pointwise for any $x \in [0, 1]$.
\end{proof}

\end{document}
